\section{SF-E842--843. A socle-pair obstruction and arbitrary parameter-power point defects}
\label{SF-E842:section}

\paragraph{Status and provenance.}
Author: \texttt{/root/sixfolds}, GPT-6 Astra Pro; inherited effort is not exposed.
The complete frozen Markdown proof was prepared on 22 November 2026 during
an execution-environment outage and preserved in runtime memory. The author,
root, and the independent \texttt{/root/g34\_eightfolds} reviewer have read the
complete proof; all three mathematical reviews passed. The peer also checked
the complete author program and produced an independent exact certificate.
The durable draft preserves its original pending-review status and all
outage-time metadata. SF-E841 remains a separate unpromoted derived draft.

\subsection{Local setup and the previously proved deformation interface}
\label{SF-E842:interface}

Let \(R=\mathbb C[[x_1,x_2,x_3,x_4]]\), let
\(I=(f_1,f_2,f_3,f_4)\) be a proper finite-colength complete-intersection
ideal, and set \(A=R/I\). Let \(\pi_+,\pi_-\) be constant rank-two
bivectors with complementary supports. Thus \(\pi_+\wedge\pi_-\) is a
nonzero scalar multiple of the coordinate volume. Write \(P_+,P_-\)
for the alternating four-by-four matrices, or equivalently their six
coefficient tuples,
\[
 P_{\pm,ij}=\pi_\pm(df_i,df_j)\pmod I,\qquad i<j.
 \tag{E842.1}
\]

Consider an actual finite \(R\)-module with the hull-modification sequence
\[
 0\longrightarrow A^n\longrightarrow M\longrightarrow R^r
       \longrightarrow A^m\longrightarrow0,\qquad m+n>0.
 \tag{E842.2}
\]

The exact E819/E836 interface is used, with no new deformation premise.
If \(M\) has flat first-order deformations for both multiplication classes
\(\pi_+\) and \(\pi_-\), the finite-perfect-cone supertrace argument gives
\(m=n>0\). The lifted map to the free hull and the ordinary Koszul extension
description then give actual square matrices
\(C_+,C_-,B_{ij}\in\operatorname{Mat}_m(A)\) with both orders
\[
 C_\pm B_{ij}=B_{ij}C_\pm=P_{\pm,ij}I_m.
 \tag{E842.3}
\]

Here the same \(B_{ij}\) are the extension matrices for both bivectors.
A chosen free summand in the hull image need not remain split under
deformation.

For precision, this interface uses
\(\operatorname{Ext}_R^j(A,R)=0\) for \(j<4\), which lifts the hull map.
Its first-order value on torsion defines \(C_\pi:T\to Q\).
Lifting the relation \(f_i e_j-f_j e_i=B_{ij}\) in \(M\) and applying that
hull map gives \(C_\pi B_{ij}=P^\pi_{ij}I\).
Applying the deformed commutator of \(f_i,f_j\) to a lifted torsion element
gives \(B_{ij}C_\pi=P^\pi_{ij}I\). These are the two complete actual-module
calculations proved in E836, Sections 2--3. No commutativity among the
\(B_{ij}\), or between \(C_+\) and \(C_-\), is assumed here.

The complete-intersection Jacobian-socle theorem is also retained with its
precise characteristic-zero hypothesis:
\[
 \Delta_f=\det(\partial f_i/\partial x_j)\pmod I
       \quad\hbox{generates the nonzero one-dimensional socle of }A.
 \tag{E842.4}
\]

The exact published source is Eisenbud--Ulrich, \emph{Duality and socle
generators for residual intersections}, \emph{J. reine angew. Math.}
756 (2019), 183--226, Theorem 7.1, p.~210, with the complete invoked proofs
in Appendix A.1, A.3 and A.5, pp.~222--225.\cite{SFEisenbudUlrich2019}
Those complete proofs were read in the earlier E836/E838/E839 reviews.
The formal-local theorem statement and the complete invoked Appendix proofs
were freshly read again during the present author and independent peer
audits. No preprint is substituted for this input.

\subsection{The elementary square-matrix socle-pair lemma}
\label{SF-E842:matrix}

\paragraph{Lemma SF-E842a.}
Let \(A\) be a commutative local complex Artinian algebra. Suppose
\(p,q\in A\) satisfy
\[
 0\ne pq\in\operatorname{Soc}(A),\qquad
                         pA\cap qA=\mathbf C\,pq.
 \tag{E842.5}
\]

If \(C,B_p,B_q\) are square matrices of positive size over \(A\) and
\[
 CB_p=B_pC=pI,\qquad CB_q=B_qC=qI,
 \tag{E842.6}
\]
then \(C\) is invertible.

\paragraph{Proof.}
The two orders, not any assumed commutation of \(B_p\) and \(B_q\), give
\[
 pB_q=B_pCB_q=qB_p.
 \tag{E842.7}
\]

Each entry of \(qB_p\) therefore belongs to \(pA\cap qA\). Consequently
\[
 qB_p=pqD
 \tag{E842.8}
\]
for a constant complex matrix \(D\). Apply \(C\) to this equality.
Since \(pq\) is a nonzero socle element, every entry of the maximal ideal
kills it, so
\[
 pqI=qCB_p=C(pqD)=pq\,C(0)D.
 \tag{E842.9}
\]

The nonzero scalar vector \(pq\) gives \(C(0)D=I\) over \(\mathbb C\).
The square complex matrix \(C(0)\) is invertible, and therefore \(C\) is
invertible over the local ring \(A\). This proves the lemma for every
matrix size and with arbitrary higher entries. No quotient deformation
or freeness under base change has been used.

\subsection{An invertible symplectic pencil companion is impossible}
\label{SF-E842:pencil}

\paragraph{Lemma SF-E842b.}
In the complete-intersection setup (E842.1)--(E842.3), no matrix
\[
                  C=aC_++bC_-,\qquad ab\ne0,
 \tag{E842.10}
\]
is invertible.

\paragraph{Proof.}
Put \(P=aP_++bP_-\). The two existing companion equations imply
\(CB_{ij}=B_{ij}C=P_{ij}I\); a deformation in a third direction is not
being assumed.

Each \(P_\pm\) is a decomposable two-form over \(A\), since it is obtained
by applying the Jacobian matrix to the decomposable constant bivector
\(\pi_\pm\). Hence its Pfaffian is zero. The Pfaffian transformation
identity gives
\[
 \operatorname{Pf}(P_\pm)=0,\qquad
 \operatorname{Pf}(P)
   =\Delta_f\,\operatorname{Pf}(a\Pi_++b\Pi_-)
   =\lambda\Delta_f,\quad \lambda\ne0.
 \tag{E842.11}
\]

Here \(\Pi\) denotes the constant skew matrix in the original \(x\)
coordinates. The coefficient \(\lambda\) is nonzero because \(ab\ne0\)
and the two planes are complementary. By (E842.4),
\(\operatorname{Pf}(P)\) is a nonzero socle element.

If \(C\) were invertible, set \(D_+=aC_+C^{-1}\),
\(D_-=bC_-C^{-1}\), and \(d_\pm=\operatorname{tr}(D_\pm)/m\in A\).
These are defined over \(A\) in characteristic zero.
Substitution of \(B_{ij}=C^{-1}P_{ij}\) gives
\[
 D_+P_{ij}=aP_{+,ij}I,\quad
 D_-P_{ij}=bP_{-,ij}I,\quad
 D_++D_-=I.
 \tag{E842.12}
\]

Taking traces yields the exact identities
\[
             aP_+=d_+P,\qquad bP_-=d_-P,\qquad d_++d_-=1.
 \tag{E842.13}
\]

Their Pfaffians give
\(d_+^2\operatorname{Pf}(P)=d_-^2\operatorname{Pf}(P)=0\).
Multiplication on the nonzero one-dimensional socle only sees the constant
term. Thus \(d_+(0)^2=d_-(0)^2=0\), so both constants are zero.
This contradicts \(d_+(0)+d_-(0)=1\). The lemma follows.
Notice that no homogeneous initial form, matrix diagonalization,
or proportionality assumption was needed.

\paragraph{Theorem SF-E842.}
Suppose that for some \(ab\ne0\), two entries \(p,q\) of
\(P=aP_++bP_-\) satisfy the socle-pair condition (E842.5).
Then no module (E842.2) has flat first-order deformations for both
\(\pi_+\) and \(\pi_-\).

\paragraph{Proof.}
If a multiplicity is zero, the already proved supertrace equality
contradicts \(m+n>0\). Otherwise use the common positive square size and
\(C=aC_++bC_-\). Lemma E842a makes \(C\) invertible, whereas E842b
forbids this. This proves the theorem. The socle-pair condition is a
sufficient, explicitly testable criterion; it is not claimed for every
complete intersection.

\subsection{All parameter-power complete intersections, with nonlinear mixed parameters}
\label{SF-E843:powers}

\paragraph{Theorem SF-E843.}
Let \(u_1,\ldots,u_4\) be any regular system of parameters of \(R\).
They may be nonlinear formal functions mixing the two Poisson planes.
For arbitrary positive integers \(a_i\), put
\[
                  I=(u_1^{a_1},u_2^{a_2},u_3^{a_3},u_4^{a_4}).
 \tag{E843.1}
\]

Then the conclusion of SF-E842 holds for every nontrivial module (E842.2).

\paragraph{Proof.}
Set \(v_i=u_i^{a_i-1}\in A\), and \(\tau=v_1v_2v_3v_4\).
The monomial basis \(0\le r_i<a_i\) proves that \(A\) has length
\(\prod_i a_i\) and socle \(\mathbb C\tau\), with \(\tau\ne0\).
Put \(E_{ij}=v_iv_j\).

Use the already existing pencil \(\pi=\pi_++\pi_-\), which is
nondegenerate. Its coefficient matrix on the cotangent vectors \(du_i\)
at the closed point is nondegenerate because those vectors form a basis.
Its nonzero Pfaffian implies that at least one complementary pair of entries
\(h_{ij}(0),h_{kl}(0)\) is nonzero, where
\(\{i,j,k,l\}=\{1,2,3,4\}\). The actual entries of \(P\) are
\[
 P_{ij}=a_i a_j E_{ij}h_{ij},\qquad
                   h_{ij}=\pi(du_i,du_j)\pmod I.
 \tag{E843.2}
\]

For that complementary pair, both \(h_{ij}\) and \(h_{kl}\) are units.
Hence the two principal ideals are \(E_{ij}A\) and \(E_{kl}A\).
In the monomial basis, the first requires maximal exponent on \(i,j\)
and the second on \(k,l\). Their intersection requires maximal exponents
on every variable and is exactly \(\mathbb C\tau\).
Their product is a nonzero constant multiple of \(\tau\).
Thus these two actual entries satisfy (E842.5), and E842 proves the assertion.

This proof handles every exponent and every nonlinear parameter change
directly in \(A\). In particular, \(A\) is never replaced by its square
quotient; no deformation, Ext group, or free module is improperly base changed.
If some \(a_i=1\), the assertion includes cases already excluded by E836.
The new scope includes all \(a_i\ge2\) in arbitrary mixed parameter coordinates.

\subsection{Exact strict cubic witnesses}
\label{SF-E843:witnesses}

Take constant Poisson coordinates \(X=(x_1,x_2,y_1,y_2)\), with
\(\pi_+=\partial_{x_1}\wedge\partial_{x_2}\) and
\(\pi_-=\partial_{y_1}\wedge\partial_{y_2}\).
Define \(w=JX\) by the integral matrix
\[
 J=\begin{pmatrix}
 1&0&1&0\\
 0&1&0&1\\
 -3&5&-9&10\\
 -8&5&-16&15
 \end{pmatrix},\qquad \det J=-20.
 \tag{E843.3}
\]

The bivector coefficient tuples in the order \(12,13,14,23,24,34\) are
\[
 p=(1,5,5,3,8,25),\qquad q=(1,10,15,9,16,25).
 \tag{E843.4}
\]

Each tuple is decomposable: its Pfaffian is respectively
\(25-40+15=0\) and \(25-160+135=0\).
Their mixed wedge coefficient is \(-20\), so the planes are transverse.
This is the earlier dense rational boundary after scaling the last two
parameters by five; the present certificate is entirely integral.

There are two explicit variants. In the first set \(u_i=w_i\).
In the second set
\[
 u_1=w_1+w_3^2,\quad u_2=w_2+w_4^2,\quad u_3=w_3,\quad u_4=w_4.
 \tag{E843.5}
\]

The second change is a polynomial automorphism with determinant one
and triangular inverse
\(w_1=u_1-u_3^2,\ w_2=u_2-u_4^2,\ w_3=u_3,\ w_4=u_4\).
In either case set \(I=(u_1^3,\ldots,u_4^3)\).
The algebra has length \(81\), embedding dimension four and socle degree
eight; \(\tau=u_1^2u_2^2u_3^2u_4^2\ne0\).

For the linear variant, all six brackets are exactly
\(P_{+,ij}=9p_{ij}E_{ij}\) and \(P_{-,ij}=9q_{ij}E_{ij}\).
In the nonlinear variant the same formula has coefficient units whose
constants are \(p_{ij},q_{ij}\).
Every scalar cross-minor vanishes in both variants:
\[
                  P_{+,e}P_{-,f}-P_{-,e}P_{+,f}=0.
 \tag{E843.6}
\]

If \(e,f\) share an index, the product contains \(u_i^4\) and is zero.
If they are complementary, the product of the \(E\) entries is \(\tau\),
which kills every positive-degree coefficient.
The relevant ratios \(q_e/p_e\) are equal on each complementary pair:
\(1\) on \(12/34\), \(2\) on \(13/24\), and \(3\) on \(14/23\).
This proves all scalar cross-minors, including all ordered pairs.

The actual generator Jacobian and the mixed bracket wedge are
\[
 \det\bigl(\partial(u_i^3)/\partial X_j\bigr)
       =-1620\,\tau,\qquad
                P_+\wedge P_-=-1620\,\tau
 \tag{E843.7}
\]
with the mixed-wedge coefficient convention of (E843.4).
Indeed \(3^4\det J=81(-20)\), and the nonlinear triangular change
has determinant one.

For \(P=P_++P_-\), choose its \(12\) and \(34\) entries.
Their constant coefficients relative to \(E_{12},E_{34}\) are \(18\)
and \(450\). Their ideals each have dimension nine, their intersection
has dimension one, and their product is \(8100\tau\).
This is the socle pair used by the new theorem.

These examples are outside the earlier local exclusions for exact reasons.
There is no linear generator and the embedding dimension is four.
The lowest ideal order is three, invariant under formal coordinate
and minimal-generator changes, so no regular quadratic initial presentation
exists; the square-thickening and quadratic-jet E837--840 families do not apply.
For each pure bivector the leading quartic bracket map from the exterior
square of the four-dimensional minimal relation space has rank six:
all six coefficients in (E843.4) are nonzero and the six \(E_{ij}\)
are independent degree-four monomials.
A triple of independent minimal generators has a three-dimensional exterior
square, which cannot lie in that zero kernel.
The exact covariance under an arbitrary invertible generator change is
the one proved in E836. Thus no minimal presentation has the commuting
triple required by E836, for either pure bivector.
Finally, (E843.6) shows that this is not an example already rejected
by the scalar companion minors.

\subsection{Actual coherent balanced examples and the equivariant consequence}
\label{SF-E843:sheaves}

Use the actual locally free balanced-character hull \(H\) from E815/E820,
or a direct sum of copies, on the explicit genuine RM abelian fourfold.
Choose a point and the finite subscheme \(Z\) with one of the preceding
completed ideals. It is an algebraic finite subscheme: the ideal contains
the ninth power of the maximal ideal, so its finite quotient is already
a quotient of the algebraic local ring modulo that power.
Thus no algebraization of an infinite formal family is asserted.

For any positive \(m\), take enough copies of \(H\) to have rank at least
\(m\), choose a surjection \(H\to\mathcal O_Z^m\) using a local basis
on this finite subscheme, and put \(K\) equal to its kernel.
There are actual extensions
\[
 0\longrightarrow O_Z^m\longrightarrow E\longrightarrow K\longrightarrow0,
 \qquad
                \operatorname{ch}(E)=\operatorname{ch}(H).
 \tag{E843.8}
\]

The torsion is precisely \(\mathcal O_Z^m\), the double dual is \(H\),
and the hull quotient is \(\mathcal O_Z^m\): \(K\) is torsion-free and
agrees with \(H\) away from codimension four.
Smoothness makes these coherent sheaves perfect, and coherence gives
ordinary negative-Ext vanishing.
Because positive Ext from \(H\) to a point-supported sheaf vanishes,
the complete-intersection Koszul resolution gives
\[
 \operatorname{Ext}^1(K,O_Z^m)\simeq
 \operatorname{Ext}^2(O_Z^m,O_Z^m),\qquad
                         \dim=6\cdot81\,m^2=486m^2.
 \tag{E843.9}
\]

Thus the examples form actual extension families, not formal characters.

The already proved E819 invariant-to-ordinary characteristic-map argument
says that weak balanced middle semiregularity for any specified coherent
effective finite translation/tensor action forces the two local
deformations for \(\pi_+,\pi_-\).
Its rational algebraic \(B\)-field version leaves these bivector projections
unchanged. Restricting at the chosen point and applying E843 therefore
excludes weak semiregularity for every extension in (E843.8).
No action on a chosen local basis, no simplicity, and no equivariant
splitting of the extension is assumed.
This does not assert a universal obstruction for unrelated Artinian modules
or arbitrary perfect complexes.

\subsection{Exact computation and check status}
\label{SF-E843:checks}

A new exact certificate ran in the available JavaScript runtime at
\texttt{2026-11-22T00:56:11.438Z}, with recorded duration \(0.006\) seconds.
The first run passed.
Every polynomial integer addition and product is guarded by
\texttt{Number.isSafeInteger}.
It reconstructs both actual gradient matrices, all \(6561\) basis products
of the length-\(81\) algebra, both pure Pfaffians, the nonzero mixed
and actual Jacobian, all \(36\) ordered scalar cross-minors in each variant,
both initial ranks, the two nine-dimensional principal ideals and their
one-dimensional socle intersection.
This finite check supplements the universal all-size proofs.

The complete runtime source and original first result were preserved
during the outage and saved after restoration as
\texttt{rm\_socle\_pair\_power\_certificate.js} and
\texttt{rm\_socle\_pair\_power\_certificate\_first\_result.json}.
The frozen v1 draft is
\texttt{rm\_socle\_pair\_power\_obstruction\_draft.md}.
A restored-runtime rerun at \texttt{2026-11-22T01:09:37.894Z} passed in
\(0.018\) seconds and had exactly the same invariant output.
The independent peer's separately authored exact reconstruction also passed;
its original failure-and-repair history concerns only its own initial
\texttt{Array.map} callback and is retained in the peer artifacts.
Root has now read the complete canonical TeX proof, the entire author
program and its first complete output, and the complete independent peer
report. The separate peer also bound its full review to this canonical
file and checked that all22 displays match the frozen draft. No fresh
root execution of the author or peer programs is claimed here.
No successful semiregular seed has been constructed.
